\section{Method}
\label{sec:method}

\begin{figure}[t]
  \centering
  \IfFileExists{figures/architecture.pdf}
    {\includegraphics[width=\columnwidth]{figures/architecture.pdf}}%
    {%% architecture-tikz.tex --- the tikzpicture body for Figure 1.
%% Single source of truth. Consumed in two ways:
%%   (a) figures/architecture.tex (standalone) for separate compile.
%%   (b) sections/04-method.tex via %% architecture-tikz.tex --- the tikzpicture body for Figure 1.
%% Single source of truth. Consumed in two ways:
%%   (a) figures/architecture.tex (standalone) for separate compile.
%%   (b) sections/04-method.tex via %% architecture-tikz.tex --- the tikzpicture body for Figure 1.
%% Single source of truth. Consumed in two ways:
%%   (a) figures/architecture.tex (standalone) for separate compile.
%%   (b) sections/04-method.tex via \input{figures/architecture-tikz}
%%       so the figure renders inline when main.tex is compiled.
%%
%% Required preamble in the host document:
%%   \usepackage{tikz}
%%   \usetikzlibrary{positioning, arrows.meta, shapes.geometric,
%%                   fit, backgrounds, calc}

\begin{tikzpicture}[
  node distance=4mm,
  >={Stealth[length=2mm]},
  every node/.style={font=\sffamily\scriptsize}
]
\tikzset{
  box/.style    ={draw=black, line width=0.5pt, rounded corners=2pt,
                  align=center, inner sep=2pt},
  light/.style  ={box, fill=black!8},
  mid/.style    ={box, fill=black!18},
  pipe/.style   ={light, text width=1.55cm, minimum width=1.7cm,
                  minimum height=7.5mm},
  lp/.style     ={light, text width=1.45cm, minimum width=1.55cm,
                  minimum height=6.5mm},
  store/.style  ={cylinder, shape border rotate=90, aspect=0.25,
                  draw=black, line width=0.5pt, fill=black!18,
                  align=center, inner sep=2pt,
                  text width=1.7cm, minimum height=7mm},
  arr/.style    ={->, line width=0.4pt},
  darr/.style   ={arr, dashed},
  lab/.style    ={font=\sffamily\tiny, fill=white, inner sep=1pt},
  group/.style  ={draw=black, line width=0.5pt, rounded corners=2pt,
                  dotted, inner sep=4pt}
}

%% ----- Top: Query, Router -----
\node[light, text width=1.7cm] (q) at (0,0) {Query $q$};
\node[mid, below=4mm of q, text width=4.4cm] (router)
  {Hop-Adaptive Router\\{\tiny predicts $\widehat{H}(q) \in \{1, 2, 3{+}\}$}};
\draw[arr] (q) -- (router);

%% ----- Pipelines (centered at x=0) -----
\node[pipe, below=5mm of router, xshift=-2.775cm] (p1)
  {Vanilla\\{\tiny embed $\to$ top-$k$ $\to$ LLM}};
\node[pipe, right=1.5mm of p1] (p2)
  {Agentic\\{\tiny LangGraph loop}};
\node[pipe, right=1.5mm of p2] (p3)
  {GraphRAG\\{\tiny seeds + 2-hop walk}};
\node[pipe, right=1.5mm of p3] (p4)
  {Agentic+Graph\\{\tiny LangGraph + graph}};
\foreach \i in {1,2,3,4}
  \draw[arr] (router.south) -- (p\i.north);

%% ----- Agentic Loop subgraph (centered below pipelines) -----
\node[lp] (critic) at (-1.7,-4.6) {Router/Critic};
\node[lp, right=1.5mm of critic] (retr)  {Retriever};
\node[lp, right=1.5mm of retr]   (tools) {Tools\\{\tiny search, graph}};
\node[lp, below=4mm of retr]     (synth) {Synthesizer};

\draw[arr] (critic) -- (retr);
\draw[arr] (retr)   -- (tools);
\draw[arr] (tools.north) to[out=90,in=90,looseness=0.5]
  node[lab, midway, yshift=1mm] {iter $\le 3$} (critic.north);
\draw[arr] (critic.south) |- (synth.west);

\begin{scope}[on background layer]
  \node[group, fit=(critic)(retr)(tools)(synth)] (loop) {};
\end{scope}
\node[lab, anchor=north west]
  at ([xshift=2pt,yshift=-1pt]loop.north west) {Agentic Loop};

%% ----- P2, P4 -> Loop (dashed "uses"; right-angle elbows) -----
\draw[darr] (p2.south) -- ++(0,-2mm) -| (loop.north west)
  node[lab, pos=0.4, above] {uses};
\draw[darr] (p4.south) -- ++(0,-2mm) -| (loop.north east)
  node[lab, pos=0.4, above] {uses};

%% ----- Output -----
\node[light, below=6mm of synth, text width=4.5cm] (out)
  {Answer $a(q)$ + cited IDs $C(q)$};
\draw[arr] (synth.south) -- (out.north);

%% ----- P1, P3 -> Output (bypass loop on the sides) -----
\draw[arr] (p1.south) -- ++(0,-2mm) -|
  ([xshift=-4mm]loop.south west) |- (out.west);
\draw[arr] (p3.south) -- ++(0,-2mm) -|
  ([xshift=4mm]loop.south east)  |- (out.east);

%% ----- Data stores at the bottom; consolidated arrows -----
\node[store, below=10mm of out, xshift=-22mm] (s1)
  {Vector Store\\{\tiny (Chroma)}};
\node[store, below=10mm of out, xshift=22mm]  (s2)
  {Graph Store\\{\tiny (Neo4j Aura)}};

\draw[darr] (s1.north) -- ++(0,4mm)
  node[lab, anchor=south] {$\to$ P1--P4};
\draw[darr] (s2.north) -- ++(0,4mm)
  node[lab, anchor=south] {$\to$ P3, P4};

\end{tikzpicture}

%%       so the figure renders inline when main.tex is compiled.
%%
%% Required preamble in the host document:
%%   \usepackage{tikz}
%%   \usetikzlibrary{positioning, arrows.meta, shapes.geometric,
%%                   fit, backgrounds, calc}

\begin{tikzpicture}[
  node distance=4mm,
  >={Stealth[length=2mm]},
  every node/.style={font=\sffamily\scriptsize}
]
\tikzset{
  box/.style    ={draw=black, line width=0.5pt, rounded corners=2pt,
                  align=center, inner sep=2pt},
  light/.style  ={box, fill=black!8},
  mid/.style    ={box, fill=black!18},
  pipe/.style   ={light, text width=1.55cm, minimum width=1.7cm,
                  minimum height=7.5mm},
  lp/.style     ={light, text width=1.45cm, minimum width=1.55cm,
                  minimum height=6.5mm},
  store/.style  ={cylinder, shape border rotate=90, aspect=0.25,
                  draw=black, line width=0.5pt, fill=black!18,
                  align=center, inner sep=2pt,
                  text width=1.7cm, minimum height=7mm},
  arr/.style    ={->, line width=0.4pt},
  darr/.style   ={arr, dashed},
  lab/.style    ={font=\sffamily\tiny, fill=white, inner sep=1pt},
  group/.style  ={draw=black, line width=0.5pt, rounded corners=2pt,
                  dotted, inner sep=4pt}
}

%% ----- Top: Query, Router -----
\node[light, text width=1.7cm] (q) at (0,0) {Query $q$};
\node[mid, below=4mm of q, text width=4.4cm] (router)
  {Hop-Adaptive Router\\{\tiny predicts $\widehat{H}(q) \in \{1, 2, 3{+}\}$}};
\draw[arr] (q) -- (router);

%% ----- Pipelines (centered at x=0) -----
\node[pipe, below=5mm of router, xshift=-2.775cm] (p1)
  {Vanilla\\{\tiny embed $\to$ top-$k$ $\to$ LLM}};
\node[pipe, right=1.5mm of p1] (p2)
  {Agentic\\{\tiny LangGraph loop}};
\node[pipe, right=1.5mm of p2] (p3)
  {GraphRAG\\{\tiny seeds + 2-hop walk}};
\node[pipe, right=1.5mm of p3] (p4)
  {Agentic+Graph\\{\tiny LangGraph + graph}};
\foreach \i in {1,2,3,4}
  \draw[arr] (router.south) -- (p\i.north);

%% ----- Agentic Loop subgraph (centered below pipelines) -----
\node[lp] (critic) at (-1.7,-4.6) {Router/Critic};
\node[lp, right=1.5mm of critic] (retr)  {Retriever};
\node[lp, right=1.5mm of retr]   (tools) {Tools\\{\tiny search, graph}};
\node[lp, below=4mm of retr]     (synth) {Synthesizer};

\draw[arr] (critic) -- (retr);
\draw[arr] (retr)   -- (tools);
\draw[arr] (tools.north) to[out=90,in=90,looseness=0.5]
  node[lab, midway, yshift=1mm] {iter $\le 3$} (critic.north);
\draw[arr] (critic.south) |- (synth.west);

\begin{scope}[on background layer]
  \node[group, fit=(critic)(retr)(tools)(synth)] (loop) {};
\end{scope}
\node[lab, anchor=north west]
  at ([xshift=2pt,yshift=-1pt]loop.north west) {Agentic Loop};

%% ----- P2, P4 -> Loop (dashed "uses"; right-angle elbows) -----
\draw[darr] (p2.south) -- ++(0,-2mm) -| (loop.north west)
  node[lab, pos=0.4, above] {uses};
\draw[darr] (p4.south) -- ++(0,-2mm) -| (loop.north east)
  node[lab, pos=0.4, above] {uses};

%% ----- Output -----
\node[light, below=6mm of synth, text width=4.5cm] (out)
  {Answer $a(q)$ + cited IDs $C(q)$};
\draw[arr] (synth.south) -- (out.north);

%% ----- P1, P3 -> Output (bypass loop on the sides) -----
\draw[arr] (p1.south) -- ++(0,-2mm) -|
  ([xshift=-4mm]loop.south west) |- (out.west);
\draw[arr] (p3.south) -- ++(0,-2mm) -|
  ([xshift=4mm]loop.south east)  |- (out.east);

%% ----- Data stores at the bottom; consolidated arrows -----
\node[store, below=10mm of out, xshift=-22mm] (s1)
  {Vector Store\\{\tiny (Chroma)}};
\node[store, below=10mm of out, xshift=22mm]  (s2)
  {Graph Store\\{\tiny (Neo4j Aura)}};

\draw[darr] (s1.north) -- ++(0,4mm)
  node[lab, anchor=south] {$\to$ P1--P4};
\draw[darr] (s2.north) -- ++(0,4mm)
  node[lab, anchor=south] {$\to$ P3, P4};

\end{tikzpicture}

%%       so the figure renders inline when main.tex is compiled.
%%
%% Required preamble in the host document:
%%   \usepackage{tikz}
%%   \usetikzlibrary{positioning, arrows.meta, shapes.geometric,
%%                   fit, backgrounds, calc}

\begin{tikzpicture}[
  node distance=4mm,
  >={Stealth[length=2mm]},
  every node/.style={font=\sffamily\scriptsize}
]
\tikzset{
  box/.style    ={draw=black, line width=0.5pt, rounded corners=2pt,
                  align=center, inner sep=2pt},
  light/.style  ={box, fill=black!8},
  mid/.style    ={box, fill=black!18},
  pipe/.style   ={light, text width=1.55cm, minimum width=1.7cm,
                  minimum height=7.5mm},
  lp/.style     ={light, text width=1.45cm, minimum width=1.55cm,
                  minimum height=6.5mm},
  store/.style  ={cylinder, shape border rotate=90, aspect=0.25,
                  draw=black, line width=0.5pt, fill=black!18,
                  align=center, inner sep=2pt,
                  text width=1.7cm, minimum height=7mm},
  arr/.style    ={->, line width=0.4pt},
  darr/.style   ={arr, dashed},
  lab/.style    ={font=\sffamily\tiny, fill=white, inner sep=1pt},
  group/.style  ={draw=black, line width=0.5pt, rounded corners=2pt,
                  dotted, inner sep=4pt}
}

%% ----- Top: Query, Router -----
\node[light, text width=1.7cm] (q) at (0,0) {Query $q$};
\node[mid, below=4mm of q, text width=4.4cm] (router)
  {Hop-Adaptive Router\\{\tiny predicts $\widehat{H}(q) \in \{1, 2, 3{+}\}$}};
\draw[arr] (q) -- (router);

%% ----- Pipelines (centered at x=0) -----
\node[pipe, below=5mm of router, xshift=-2.775cm] (p1)
  {Vanilla\\{\tiny embed $\to$ top-$k$ $\to$ LLM}};
\node[pipe, right=1.5mm of p1] (p2)
  {Agentic\\{\tiny LangGraph loop}};
\node[pipe, right=1.5mm of p2] (p3)
  {GraphRAG\\{\tiny seeds + 2-hop walk}};
\node[pipe, right=1.5mm of p3] (p4)
  {Agentic+Graph\\{\tiny LangGraph + graph}};
\foreach \i in {1,2,3,4}
  \draw[arr] (router.south) -- (p\i.north);

%% ----- Agentic Loop subgraph (centered below pipelines) -----
\node[lp] (critic) at (-1.7,-4.6) {Router/Critic};
\node[lp, right=1.5mm of critic] (retr)  {Retriever};
\node[lp, right=1.5mm of retr]   (tools) {Tools\\{\tiny search, graph}};
\node[lp, below=4mm of retr]     (synth) {Synthesizer};

\draw[arr] (critic) -- (retr);
\draw[arr] (retr)   -- (tools);
\draw[arr] (tools.north) to[out=90,in=90,looseness=0.5]
  node[lab, midway, yshift=1mm] {iter $\le 3$} (critic.north);
\draw[arr] (critic.south) |- (synth.west);

\begin{scope}[on background layer]
  \node[group, fit=(critic)(retr)(tools)(synth)] (loop) {};
\end{scope}
\node[lab, anchor=north west]
  at ([xshift=2pt,yshift=-1pt]loop.north west) {Agentic Loop};

%% ----- P2, P4 -> Loop (dashed "uses"; right-angle elbows) -----
\draw[darr] (p2.south) -- ++(0,-2mm) -| (loop.north west)
  node[lab, pos=0.4, above] {uses};
\draw[darr] (p4.south) -- ++(0,-2mm) -| (loop.north east)
  node[lab, pos=0.4, above] {uses};

%% ----- Output -----
\node[light, below=6mm of synth, text width=4.5cm] (out)
  {Answer $a(q)$ + cited IDs $C(q)$};
\draw[arr] (synth.south) -- (out.north);

%% ----- P1, P3 -> Output (bypass loop on the sides) -----
\draw[arr] (p1.south) -- ++(0,-2mm) -|
  ([xshift=-4mm]loop.south west) |- (out.west);
\draw[arr] (p3.south) -- ++(0,-2mm) -|
  ([xshift=4mm]loop.south east)  |- (out.east);

%% ----- Data stores at the bottom; consolidated arrows -----
\node[store, below=10mm of out, xshift=-22mm] (s1)
  {Vector Store\\{\tiny (Chroma)}};
\node[store, below=10mm of out, xshift=22mm]  (s2)
  {Graph Store\\{\tiny (Neo4j Aura)}};

\draw[darr] (s1.north) -- ++(0,4mm)
  node[lab, anchor=south] {$\to$ P1--P4};
\draw[darr] (s2.north) -- ++(0,4mm)
  node[lab, anchor=south] {$\to$ P3, P4};

\end{tikzpicture}
}%
  \caption{Architecture. A hop-adaptive router selects among four
    retrieval pipelines based on a query-conditioned hop-distance
    estimator. Pipelines~2 and~4 share a LangGraph
    router--retriever--critic--synthesizer skeleton; the typed
    graph-lookup tool is exposed in pipelines~3 and~4. Tools
    abbreviated: \texttt{search} = \texttt{search\_documents},
    \texttt{graph} = \texttt{graph\_lookup}.}
  \label{fig:arch}
\end{figure}

\subsection{Problem Formulation}
\label{sec:problem}
We model the requirement corpus as a typed graph $G=(V,E,\tau)$ where
$V$ is the set of requirements, $E\subseteq V\times V$ is the set of
traceability links, and $\tau:E\to T$ assigns each edge a type from
$T=\{\mathrm{derives\_from},\mathrm{satisfies},\mathrm{references},
\mathrm{traces\_to}\}$. A traceability query $q$ has an unobserved hop
class $H(q)\in\{1,2,3{+}\}$ that summarizes the minimum graph distance
between the entities mentioned in $q$ and the gold answer. The system
produces an answer $a(q)$ together with cited requirement IDs
$C(q)\subseteq V$.

\subsection{Four Retrieval Pipelines}
\label{sec:pipelines}
We consider four pipelines that share an embedder, a vector store, and
a synthesizer prompt and differ only in retrieval strategy:
\textbf{(i)~Vanilla} performs dense top-$k$ retrieval over text chunks
followed by one synthesis call.
\textbf{(ii)~Agentic} runs a LangGraph
router--retriever--critic--synthesizer loop with a
\texttt{search\_documents} tool and an iteration cap of three.
\textbf{(iii)~GraphRAG} retrieves vector seeds and then walks the typed
graph for up to two hops via a Cypher query against
Neo4j, synthesizing once at the end.
\textbf{(iv)~Agentic+Graph} extends pipeline~(ii) with a
\texttt{graph\_lookup} tool node alongside \texttt{search\_documents},
allowing the critic to request typed-edge expansion when text retrieval
fails.

\subsection{Hop-Adaptive Router}
\label{sec:router}
The router predicts $\widehat{H}(q)\in\{1,2,3{+}\}$ from $q$ alone and
selects the pipeline with the best expected reward in that hop bucket
on a held-out validation set. Features include (a)~entity count via
named-entity recognition, (b)~explicit requirement-ID patterns
(regex \texttt{REQ-[A-Z]+-\textbackslash d+}), (c)~question-type cues
(e.g., \textit{trace from}, \textit{how does $X$ relate to $Y$}), and
(d)~the embedding-space norm of $q$. We instantiate two variants:
\textbf{V1}, a transparent rule-based router, and \textbf{V2}, a
multinomial logistic regressor trained on 30\% of labelled queries with
a held-out test split. Hop labels for training are obtained by walking
$G$ from gold answer chains at evaluation-set construction time.

\subsection{Link-Type-Aware Graph Tool}
\label{sec:graphtool}
The graph-lookup tool has signature
\[
  \texttt{graph\_lookup}(\mathit{seed\_ids},\,\mathit{hops},\,
                         \mathit{edge\_types})
\]
and returns chunks ranked by a hybrid score
\[
  s(c\mid q) \;=\;
    \lambda_{\mathrm{text}}\,\cos(q,c)
    \;+\; \lambda_{\mathrm{graph}}\,
          \frac{w(\tau(e_{c}))}{\mathrm{dist}(\mathit{seed},c)},
\]
where $e_c$ is the edge that brought $c$ into the candidate set,
$w:T\to\mathbb{R}_+$ assigns DO-178C-aligned semantic priorities
($\mathrm{derives\_from}>\mathrm{satisfies}>
\mathrm{references}>\mathrm{traces\_to}$), and $\lambda_{\mathrm{text}},
\lambda_{\mathrm{graph}}$ are tuned on a validation split. The tool
filters out edge types not in $T$ to avoid retrieval dilution from
auxiliary schema relations such as \textit{ALLOCATED\_TO} or
\textit{IMPLEMENTS}.
